\section{Reproducibility}\label{sec:repro}

A measurement study is only useful if someone else can repeat it. Everything in this paper comes out of the artifact
repository through documented commands.
\begin{itemize}
\item \texttt{make audit} reruns the attribution audit, with the sources fetched at their pins outside the repository.
\item \texttt{make static} measures the byte accounting and the Hub metadata again from public files.
\texttt{make test} runs the test suite, including the guard against numbers without data described below, and
\texttt{make analyze figures paper} regenerates every table, figure and value and builds this PDF.
\item \texttt{make stack-fetch} and \texttt{make pod-bootstrap VARIANT=\ldots} build each variant from pinned inputs
with recorded hashes. \texttt{make bench-E\textit{n}} and \texttt{make quality-E\textit{n}} run each experiment on a
prepared pod, and \texttt{make plan} regenerates the plan of experiments.
\end{itemize}

\paragraph{No number without a file.}
Every measured number in the text is printed by a macro that looks up a value written by the analysis pipeline. The
pipeline refuses values without an existing source file, refuses controlled values computed from exploratory files,
and records the sha256 of every source (\datafile{paper/generated/provenance.json}). Tables appear only for experiments
with data, so running one of the experiments of Section~\ref{sec:not-run} (\texttt{make bench-E\textit{n}}) adds its
table on the next build, and \texttt{STRICT=1} makes the build fail if any cited value lacks its source. A test fails
if a decimal number appears in the paper sources outside these macros, unless it is a whitelisted configuration
constant with its source (\datafile{paper/literal_numbers.yaml}).

\paragraph{From a finding to its data.}
Table~\ref{tab:trace} follows each row of Table~\ref{tab:key-results} back to the experiment, the script that ran it on
the pod, the folder that holds its raw data and the table or figure that prints it. The runners are the session
scripts of Appendix~\ref{app:plan}, and each folder has a README with the status of every run in it.

\begin{table}[h]
\centering\footnotesize
\begin{tabular}{>{\raggedright\arraybackslash}p{3.2cm}l>{\raggedright\arraybackslash}p{3.4cm}>{\raggedright\arraybackslash}p{3.6cm}>{\raggedright\arraybackslash}p{2.9cm}}
\toprule
finding & exp. & ran with & folder in \texttt{data/controlled/} & printed in \\
\midrule
fits only with the offload & E0, E2 & \texttt{make static}, \repolink{experiments/runs/lean.sh}{\texttt{lean.sh}} & \repolink{data/controlled/E2-feasibility/}{\texttt{E2-feasibility/}} (bytes: \datafile{data/static/weights/}) & Tables~\ref{tab:static-placement}, \ref{tab:ctl-E2} \\
reading against writing & E12, E3 & \repolink{experiments/runs/long.sh}{\texttt{long.sh}} & \repolink{data/controlled/E12-long-context/}{\texttt{E12-long-context/}} & Figure~\ref{fig:read-write}, Table~\ref{tab:ctl-E12} \\
one user and many & E8 & \repolink{experiments/runs/lean.sh}{\texttt{lean.sh}} & \repolink{data/controlled/E8-main/}{\texttt{E8-main/}} & Figure~\ref{fig:ctl-E8}, Table~\ref{tab:ctl-E8} \\
quality against the API & E10, E11 & \repolink{experiments/runs/lean.sh}{\texttt{lean.sh}}, \repolink{quality/api_proxy.py}{\texttt{api\_proxy.py}} & \repolink{data/controlled/E10-quality-exl3/}{\texttt{E10-quality-exl3/}}, \repolink{data/controlled/E11-quality-reference/}{\texttt{E11-quality-reference/}} & Table~\ref{tab:ctl-quality-paired} \\
streams and developers & E13 & \repolink{experiments/runs/agent.sh}{\texttt{agent.sh}} & \repolink{data/controlled/E13-coding-agent/}{\texttt{E13-coding-agent/}} & Figures~\ref{fig:agent}, \ref{fig:two-capacities}, Tables~\ref{tab:agent-capacity}, \ref{tab:agent-seats} \\
upstream offload alone & E3 & \repolink{experiments/runs/lean.sh}{\texttt{lean.sh}}, \repolink{experiments/runs/long.sh}{\texttt{long.sh}} & \repolink{data/controlled/E3-host-share/}{\texttt{E3-host-share/}} & Table~\ref{tab:ctl-E3} \\
the host matters & E3 & \repolink{experiments/runs/lean.sh}{\texttt{lean.sh}}, \repolink{experiments/runs/long.sh}{\texttt{long.sh}} & \repolink{data/controlled/E3-host-share/}{\texttt{E3-host-share/}}, \repolink{data/controlled/LEAN/}{\texttt{LEAN/}} & Tables~\ref{tab:cross-pod}, \ref{tab:pod-hardware} \\
\bottomrule
\end{tabular}
\caption{Traceability of the findings in Table~\ref{tab:key-results}. Session runners are in
\datafile{experiments/runs/}, the API proxy in \datafile{quality/}.}
\label{tab:trace}
\end{table}

\paragraph{Licenses and redistribution.}
The repository contains no third-party source code. The recipe, the plugin and the kernels are fetched and patched at
build time from pinned revisions. vllm-exl3 and Diffbot's patch to it are AGPL-3.0, so a built image that includes
them carries AGPL obligations when distributed or served with modifications. Our code is Apache-2.0, compatible with
vLLM's license and combinable with the AGPL components in a built image, and text and data are CC-BY-4.0. Model weights
are never redistributed (\datafile{THIRD_PARTY.md}, \datafile{PROVENANCE.md}).

\paragraph{Archival.}
The repository is prepared for an archived release with a persistent identifier (\datafile{CITATION.cff},
\datafile{.zenodo.json}, \datafile{RELEASE_CHECKLIST.md}).

\paragraph{Availability.}
The artifact repository, \url{https://github.com/inferencekey/vllm-exl3-moe-hostoffload}, holds the code, the
stack configurations, the workloads, the raw data and the sources of this paper, and every file named in the text
links to it. The configurations that
make the model boot (the variant builds, the expert plans and the session runners) are the ones measured here and
can be reused to serve the model on the same class of hardware.
